\documentclass[11pt]{article}
\usepackage[margin=1in]{geometry}
\usepackage{amsmath,amssymb,amsthm}
\usepackage{xcolor}
\newtheorem{lemma}{Lemma}
\newtheorem*{remark}{Remark}
\newcommand{\Var}{\operatorname{Var}}
\newcommand{\edit}[1]{\par\medskip\noindent\colorbox{gray!12}{\parbox{\dimexpr\linewidth-2\fboxsep}{\small\textbf{#1}}}\par\smallskip}
\setlength{\parindent}{0pt}
\setlength{\parskip}{4pt}

\begin{document}
\begin{center}\large\bf Section 5.4: edits for the portfolio variance bound\end{center}

The variance penalty weight $(p_i-s_i)^2\sigma_i^2$ is replaced by
$(p_i-s_i)^2\sigma_i^2/\mu_i^2$. All other text of Section~5.4 is unchanged.

\edit{1. Replace the text from ``For the distribution-free setting, the following
proposition also establishes \ldots'' (before Proposition 7) through Proposition 7,
its proof, and the paragraph ``We propose a distribution-free mean-variance
newsvendor problem \ldots\ for each $i\in N$.'' by:}

Expression (18) depends on the entire distribution $F_i$. The following bound
depends only on $(\mu_i,\sigma_i)$, as does $\theta_i$, and grows with the order
quantity.

\begin{lemma}\label{lem:cv}
Let $D\ge0$ have mean $\mu>0$ and variance $\sigma^2$. Then
$\Var(\min\{D,x\})\le(\sigma/\mu)^2x^2$ for all $x\ge0$, and the constant
$(\sigma/\mu)^2$ cannot be improved.
\end{lemma}

\begin{proof}
Let $Y=\min\{D,x\}$. We use that for every random variable $Y$ with finite second
moment and every constant $c\in\mathbb R$,
\[
\Var(Y)=E[(Y-E[Y])^2]\le E[(Y-c)^2], \tag{a}
\]
since $E[(Y-c)^2]=\Var(Y)+(E[Y]-c)^2$.

\emph{Case $x\ge\mu$.} Take $c=\mu$ in (a). If $D\le x$ then $Y-\mu=D-\mu$; if
$D>x$ then $0\le x-\mu=Y-\mu<D-\mu$. Hence $|Y-\mu|\le|D-\mu|$ and
\[
\Var(Y)\le E[(Y-\mu)^2]\le E[(D-\mu)^2]=\sigma^2\le\frac{\sigma^2}{\mu^2}x^2 .
\]

\emph{Case $x<\mu$.} Take $c=x$ in (a). Then $(Y-x)^2=((x-D)^+)^2$, where
$t^+=\max\{t,0\}$. If $D<x$ then, since $D\ge0$ and $x/\mu<1$,
$0<x-D\le x-\tfrac{x}{\mu}D=\tfrac{x}{\mu}(\mu-D)$; if $D\ge x$ then $(x-D)^+=0$.
Hence $((x-D)^+)^2\le\tfrac{x^2}{\mu^2}(\mu-D)^2$ for every outcome, and
\[
\Var(Y)\le E\big[((x-D)^+)^2\big]\le\frac{x^2}{\mu^2}E[(D-\mu)^2]=\frac{\sigma^2}{\mu^2}x^2 .
\]

\emph{Sharpness.} For $q\in(0,1)$ let $P(D=0)=q$ and $P(D=\mu/(1-q))=1-q$, so
$E[D]=\mu$, and let $x=\mu/(1-q)$. Then $Y=D$, so $\Var(Y)=\sigma^2$, while the
bound equals $\sigma^2/(1-q)^2$. The ratio $(1-q)^2$ tends to $1$ as $q\to0$, so no
constant smaller than $\sigma^2/\mu^2$ is valid for all distributions.
\end{proof}

\begin{remark}
Lemma~\ref{lem:cv} also follows from the Lorenz order: for $g\ge0$ nondecreasing
with $g(t)/t$ nonincreasing, $g(D)$ is Lorenz-smaller than $D$ (Fellman, 1976;
Arnold, 1987). With $g(t)=\min\{t,x\}$, the coefficient of variation of sales does
not exceed $\sigma/\mu$, so
$\Var(\min\{D,x\})\le(\sigma/\mu)^2\min\{x,\mu\}^2$; we use the convex bound
$(\sigma/\mu)^2x^2$.
\end{remark}

Since $V[\pi_i(x_i)]=(p_i-s_i)^2\Var(\min\{D_i,x_i\})$, for independent demands
Lemma~\ref{lem:cv} gives
\[
V\Big[\sum_{i\in N}\pi_i(x_i)\Big]\le\sum_{i\in N}(p_i-s_i)^2\frac{\sigma_i^2}{\mu_i^2}x_i^2
=\Var\Big(\sum_{i\in N}(p_i-s_i)\,x_i\frac{D_i}{\mu_i}\Big), \tag{18$'$}
\]
the Markowitz variance of a portfolio with positions $(p_i-s_i)x_i$ in independent
unit-mean assets $D_i/\mu_i$. (For correlated demands, Minkowski's inequality gives
the bound $(\sum_i(p_i-s_i)(\sigma_i/\mu_i)x_i)^2$ instead.)
We propose a distribution-free mean-variance model that adds $\lambda>0$ times the
bound (18$'$) to the Gallego--Moon cost bound. Each product's cost
$\theta_i(x_i)+\lambda(p_i-s_i)^2(\sigma_i/\mu_i)^2x_i^2$ is strictly convex, and
$x_i^U$ denotes its unconstrained minimizer, the root of a quartic first-order
condition.

\edit{2. In (19a) and (20) replace $(p_i-s_i)^2\sigma_i^2$ by
$(p_i-s_i)^2\sigma_i^2/\mu_i^2$:}
\[
\text{(19a)}\quad \min_{x,\theta,z}\ \sum_{i\in N}\theta_i+\lambda\sum_{i\in N}(p_i-s_i)^2\frac{\sigma_i^2}{\mu_i^2}x_i^2,
\qquad
\text{(20)}\quad \ldots\sum_{i\in N}\theta_i+\lambda\sum_{i\in N}(p_i-s_i)^2\frac{\sigma_i^2}{\mu_i^2}u_i\ \ldots
\]

\edit{3. In ``A data-dependent range for the penalty parameter'' update $mc_i$ and
$\lambda_{\max}$, and add the last sentence:}

$mc_i(x_i)=\partial\theta_i/\partial x_i(x_i)+2\lambda(p_i-s_i)^2(\sigma_i/\mu_i)^2x_i$, and
\[
\lambda_{\max}=\max_{i\in N:\,\sigma_i>0}\Big[\frac{p_i-s_i}{2}\Big(1+\frac{\mu_i-\varepsilon}{\sqrt{\sigma_i^2+(\varepsilon-\mu_i)^2}}\Big)-(c_i-s_i)\Big]_+\frac{\mu_i^2}{2(p_i-s_i)^2\sigma_i^2\,\varepsilon}.
\]
Because the weights $(\sigma_i/\mu_i)^2$ are dimensionless, $\lambda$ and
$\lambda_{\max}$ do not depend on the unit of demand (with prices per unit).

\end{document}
